%% kapa: exam=1 round=31 year=2020 period=1 form=A subject=civil_law subject_ko=민법 seq=14 number=14
%% source: Q-Net 감정평가사 기출문제 (cst00309) artlSeq=5207219 file=extracted/1008/5207219/A형/A형/2020년도 제31회 감정평가사 1차 시험_1교시_A형.pdf
%% conversion: kordoc 4.24.51 -> tools/gen1.py (auto); review=proofread-sonnet; flags=none
\begin{question}{14}
의사표시의 효력발생시기에 관한 설명으로 옳지 않은 것은? (다툼이 있으면 판례에 따름)
\begin{choices}
  \item 상대방 있는 의사표시는 상대방에게 도달한 때에 그 효력이 생긴다.
  \item 표의자가 의사표시의 통지를 발송한 후 제한능력자가 되어도 그 의사표시의 효력은 영향을 받지 아니한다.
  \item 상대방이 현실적으로 통지를 수령하거나 그 내용을 안 때에 도달한 것으로 본다.
  \item 상대방이 정당한 사유 없이 통지의 수령을 거절한 경우, 상대방이 그 통지의 내용을 알 수 있는 객관적 상태에 놓여 있는 때에 의사표시의 효력이 생긴다.
  \item 등기우편으로 발송된 경우, 상당한 기간 내에 도달하였다고 추정된다.
\end{choices}
\end{question}
\answer{3}
